% GENERATED FILE. Edit canonical lessons, then run npm run generate:books.
% canonical-source-hash: fnv1a64:db8912b11fedd3af
% canonical-lessons: MR-C361-mahsul, MR-C361-arthasankalpa, MR-C361-nidhi, MR-C361-anudan, MR-C361-arthik

\chapter{Money and the State}
\label{ch:mr-money-and-the-state}
\hlchaptermodality{361}

\begin{chapteropening}
\textbf{By the end of this chapter:} \emph{I can talk about revenue, budgets, funds and grants.}

Complete the last lesson of chapter 361: I can talk about revenue, budgets, funds and grants.
\end{chapteropening}

\section[\texorpdfstring{\mr{महसूल} (mahsūl)}{महसूल (mahsūl)}]{\mr{महसूल} (mahsūl) — revenue; land tax}

\label{lesson:MR-C361-mahsul}



\begin{quote}
Before the new one: say the Marathi for poverty, then the Marathi for literacy.
\end{quote}

\subsection*{You'll want to know: \mr{महसूल}}
\textbf{\mr{महसूल}} — \emph{mahsūl} — \textquotedblleft{}revenue; land tax\textquotedblright{}.

\begin{quote}
\mr{सरकारला} \mr{करातून} \mr{मोठा} \mr{महसूल} \mr{मिळतो}.
\end{quote}

\begin{quote}
\emph{sarkārlā karātūn moṭhā mahsūl miḷto.}
\end{quote}

\begin{quote}
The government gets a large revenue from taxes.
\end{quote}

Masculine.

Masculine: \mr{मोठा} \mr{महसूल}; mostly money the state takes in, from taxes and land.

\begin{cousinweb}
\textbf{\mr{महसूल}} (\emph{mahsūl}) is a loan from Persian, usually traced to Arabic \emph{maḥṣūl}, \textquotedblleft{}produce, yield\textquotedblright{}: what the land yields to the state.
\end{cousinweb}

\begin{grammarlens}[title={What you've built}]
One more word for giving reasons and opinions.
\end{grammarlens}

\subsection*{Your turn}
\begin{itemize}
\raggedright
  \item \emph{Say it:} \emph{mahsūl}
  \item \emph{Say it:} \emph{mahsūl}, once more
  \item \emph{Say it:} \emph{sarkārlā karātūn moṭhā mahsūl miḷto.} — The government gets a large revenue from taxes.
  \item \emph{Recall:} say the Marathi for poverty, then the Marathi for literacy, then say \emph{mahsūl} again
\end{itemize}

\begin{tcolorbox}[breakable,colback=teal!4,colframe=teal!35!black,title={Before you move on}]
What does \mr{महसूल} mean? (\textquotedblleft{}Revenue; land tax\textquotedblright{}.) Say it once more.
\end{tcolorbox}
\section[\texorpdfstring{\mr{अर्थसंकल्प} (arthasaṅkalpa)}{अर्थसंकल्प (arthasaṅkalpa)}]{\mr{अर्थसंकल्प} (arthasaṅkalpa) — a budget}

\label{lesson:MR-C361-arthasankalpa}



\begin{quote}
Before the new one: say the Marathi for literacy, then the Marathi for revenue.
\end{quote}

\subsection*{You'll want to know: \mr{अर्थसंकल्प}}
\textbf{\mr{अर्थसंकल्प}} — \emph{arthasaṅkalpa} — \textquotedblleft{}a budget\textquotedblright{}.

\begin{quote}
\mr{सरकारने} \mr{या} \mr{वर्षीचा} \mr{अर्थसंकल्प} \mr{मांडला}.
\end{quote}

\begin{quote}
\emph{sarkārne yā varṣīcā arthasaṅkalpa māṇḍlā.}
\end{quote}

\begin{quote}
The government presented this year's budget.
\end{quote}

Masculine.

\mr{अर्थसंकल्प} \mr{मांडणे}, to present a budget.

\begin{cousinweb}
\textbf{\mr{अर्थसंकल्प}} (\emph{arthasaṅkalpa}) is a learned compound built from Sanskrit \emph{artha}, \textquotedblleft{}wealth, aim\textquotedblright{}, and \emph{saṁkalpa}, \textquotedblleft{}will, resolve\textquotedblright{} (Monier-Williams): a plan for money.
\end{cousinweb}

\begin{grammarlens}[title={What you've built}]
One more word for giving reasons and opinions.
\end{grammarlens}

\subsection*{Your turn}
\begin{itemize}
\raggedright
  \item \emph{Say it:} \emph{arthasaṅkalpa}
  \item \emph{Say it:} \emph{arthasaṅkalpa}, once more
  \item \emph{Say it:} \emph{sarkārne yā varṣīcā arthasaṅkalpa māṇḍlā.} — The government presented this year's budget.
  \item \emph{Recall:} say the Marathi for literacy, then the Marathi for revenue, then say \emph{arthasaṅkalpa} again
\end{itemize}

\begin{tcolorbox}[breakable,colback=teal!4,colframe=teal!35!black,title={Before you move on}]
What does \mr{अर्थसंकल्प} mean? (\textquotedblleft{}A budget\textquotedblright{}.) Say it once more.
\end{tcolorbox}
\section[\texorpdfstring{\mr{निधी} (nidhī)}{निधी (nidhī)}]{\mr{निधी} (nidhī) — a fund}

\label{lesson:MR-C361-nidhi}



\begin{quote}
Before the new one: say the Marathi for revenue, then the Marathi for a budget.
\end{quote}

\subsection*{You'll want to know: \mr{निधी}}
\textbf{\mr{निधी}} — \emph{nidhī} — \textquotedblleft{}a fund\textquotedblright{}.

\begin{quote}
\mr{आम्ही} \mr{ग्रंथालयासाठी} \mr{निधी} \mr{गोळा} \mr{केला}.
\end{quote}

\begin{quote}
\emph{āmhī granthālyāsāṭhī nidhī goḷā kelā.}
\end{quote}

\begin{quote}
We raised a fund for the library.
\end{quote}

Masculine.

Masculine: \mr{मोठा} \mr{निधी}.

\begin{cousinweb}
\textbf{\mr{निधी}} (\emph{nidhī}) is a tatsama, taken from Sanskrit \emph{nidhi}, \textquotedblleft{}setting down, a hoard\textquotedblright{} (CDIAL 7207).
\end{cousinweb}

\begin{grammarlens}[title={What you've built}]
One more word for giving reasons and opinions.
\end{grammarlens}

\subsection*{Your turn}
\begin{itemize}
\raggedright
  \item \emph{Say it:} \emph{nidhī}
  \item \emph{Say it:} \emph{nidhī}, once more
  \item \emph{Say it:} \emph{āmhī granthālyāsāṭhī nidhī goḷā kelā.} — We raised a fund for the library.
  \item \emph{Recall:} say the Marathi for revenue, then the Marathi for a budget, then say \emph{nidhī} again
\end{itemize}

\begin{tcolorbox}[breakable,colback=teal!4,colframe=teal!35!black,title={Before you move on}]
What does \mr{निधी} mean? (\textquotedblleft{}A fund\textquotedblright{}.) Say it once more.
\end{tcolorbox}
\section[\texorpdfstring{\mr{अनुदान} (anudān)}{अनुदान (anudān)}]{\mr{अनुदान} (anudān) — a subsidy; a grant}

\label{lesson:MR-C361-anudan}



\begin{quote}
Before the new one: say the Marathi for a budget, then the Marathi for a fund.
\end{quote}

\subsection*{You'll want to know: \mr{अनुदान}}
\textbf{\mr{अनुदान}} — \emph{anudān} — \textquotedblleft{}a subsidy; a grant\textquotedblright{}.

\begin{quote}
\mr{या} \mr{शाळेला} \mr{सरकारकडून} \mr{अनुदान} \mr{मिळतं}.
\end{quote}

\begin{quote}
\emph{yā śāḷelā sarkārkaḍūn anudān miḷtaṁ.}
\end{quote}

\begin{quote}
This school gets a grant from the government.
\end{quote}

Neuter.

Neuter; \mr{अनुदानित}, aided (school).

\begin{cousinweb}
\textbf{\mr{अनुदान}} (\emph{anudān}) is a learned word built from Sanskrit \emph{anu-}, \textquotedblleft{}along, after\textquotedblright{}, and \emph{dāna}, \textquotedblleft{}giving\textquotedblright{} (Monier-Williams).
\end{cousinweb}

\begin{grammarlens}[title={What you've built}]
One more word for giving reasons and opinions.
\end{grammarlens}

\subsection*{Your turn}
\begin{itemize}
\raggedright
  \item \emph{Say it:} \emph{anudān}
  \item \emph{Say it:} \emph{anudān}, once more
  \item \emph{Say it:} \emph{yā śāḷelā sarkārkaḍūn anudān miḷtaṁ.} — This school gets a grant from the government.
  \item \emph{Recall:} say the Marathi for a budget, then the Marathi for a fund, then say \emph{anudān} again
\end{itemize}

\begin{tcolorbox}[breakable,colback=teal!4,colframe=teal!35!black,title={Before you move on}]
What does \mr{अनुदान} mean? (\textquotedblleft{}A subsidy; a grant\textquotedblright{}.) Say it once more.
\end{tcolorbox}
\section[\texorpdfstring{\mr{आर्थिक} (ārthik)}{आर्थिक (ārthik)}]{\mr{आर्थिक} (ārthik) — economic; financial}

\label{lesson:MR-C361-arthik}



\begin{quote}
Before the new one: say the Marathi for a fund, then the Marathi for a subsidy.
\end{quote}

\subsection*{You'll want to know: \mr{आर्थिक}}
\textbf{\mr{आर्थिक}} — \emph{ārthik} — \textquotedblleft{}economic; financial\textquotedblright{}.

\begin{quote}
\mr{त्यांना} \mr{सध्या} \mr{आर्थिक} \mr{अडचण} \mr{आहे}.
\end{quote}

\begin{quote}
\emph{tyānnā sadhyā ārthik aḍcaṇ āhe.}
\end{quote}

\begin{quote}
They are in financial difficulty at the moment.
\end{quote}

Invariable; built on \mr{अर्थ}, wealth.

\begin{cousinweb}
\textbf{\mr{आर्थिक}} (\emph{ārthik}) is a tatsama, Sanskrit \emph{ārthika} (Monier-Williams), built on \emph{artha}, \textquotedblleft{}wealth, aim\textquotedblright{}; Marathi uses it for money matters.
\end{cousinweb}

\begin{grammarlens}[title={What you've built}]
One more word for giving reasons and opinions.
\end{grammarlens}

\subsection*{Your turn}
\begin{itemize}
\raggedright
  \item \emph{Say it:} \emph{ārthik}
  \item \emph{Say it:} \emph{ārthik}, once more
  \item \emph{Say it:} \emph{tyānnā sadhyā ārthik aḍcaṇ āhe.} — They are in financial difficulty at the moment.
  \item \emph{Recall:} say the Marathi for a fund, then the Marathi for a subsidy, then say \emph{ārthik} again
\end{itemize}

\begin{tcolorbox}[breakable,colback=teal!4,colframe=teal!35!black,title={Before you move on}]
What does \mr{आर्थिक} mean? (\textquotedblleft{}Economic; financial\textquotedblright{}.) Say it once more.
\end{tcolorbox}
